\section{通过威尔逊线重正化改进准夸克分布}

让我们先回顾准夸克分布的定义。对于非极化夸克密度，其定义为~\cite{Ji:2013dva}
\beq\label{qUnpolDef}
   \tilde q(x, \Lambda, p^z) = \int^\infty_{-\infty} \frac{dz}{4\pi} e^{izk^z}  \langle p|\overline{\psi}(0, 0_\perp, z)
   \gamma^z L(z,0)\psi(0) |p\rangle \ ,
\eeq
其中夸克场沿空间 $z$ 方向分开，$L(z,0)$ 是为保证规范不变性而插入的威尔逊线规范连接，$\Lambda$ 表示紫外截断。

上述夸克密度的单圈修正既在轴规范~\cite{Xiong:2013bka} 中、也在费曼规范~\cite{Ji:2015jwa} 中被计算过。在文献~\cite{Ji:2015jwa} 中，威尔逊线 $L(z,0)$ 被拆分为沿 $z$ 轴的两段威尔逊线 $L(\infty,0)$ 与 $L(z,\infty)$，并分别与夸克场组合成规范不变的结合体。此处我们保持规范连接为 $L(z,0)$。单圈图由图~\ref{1loopFeyn} 给出。容易验证，这些图给出的单圈结果与文献~\cite{Ji:2015jwa} 中的相同。线性发散来自最后一幅图，即威尔逊线自能图。众所周知，借助辅助 $z$ 场形式~\cite{Dotsenko:1979wb,Dorn:1986dt,Craigie:1980qs}，这种线性发散可以通过质量重正化予以消除；在该形式中，$z$ 场定义在一维参数空间上，非定域的威尔逊线可诠释为 $z$ 场的两点函数。在某种意义上，辅助场 $Z(z)$ 可理解为延伸于 $[z,\infty]$ 之间的威尔逊线，即
\beq
Z(z)=L(z,\infty),
\eeq
它满足运动方程
\beq
\left[\partial_z-i g A_z(z)\right]Z(z)=0 ,
\eeq
这类似于（重夸克极限下的）重夸克场。而威尔逊线
\beq
L(z,0)=Z(z)Z^{\dagger}(0)
\eeq
的重正化方式也类似于重夸克两点函数：
\beq\label{WLrenorm}
L^{\text{ren}}(z,0)=\mathcal{Z}_Z^{-1}e^{-\delta m |z|}L(z,0) ,
\eeq
其中 $\delta m$ 类似于量纲为一的重夸克质量反项且线性发散，而量纲为零的波函数重正化常数 $\mathcal{Z}_Z$ 对数发散~\cite{Dotsenko:1979wb}。

\begin{figure}[tbp]
\centering
\includegraphics[width=0.2\textwidth]{images/1loop1}
\hspace{2em}
\includegraphics[width=0.2\textwidth]{images/1loop2}
\hspace{2em}
\includegraphics[width=0.2\textwidth]{images/1loop3}
\hspace{2em}
\includegraphics[width=0.2\textwidth]{images/1loop4}
\caption{费曼规范下准夸克分布的单圈图，共轭图未画出。}
\label{1loopFeyn}
%\vspace*{-.5em}
\end{figure}

为了展示来自质量反项的线性发散如何抵消，让我们在单圈层次确定 $\delta m$。回顾图~\ref{1loopFeyn} 的最后一幅图给出如下线性发散
\beq
\lim_{\epsilon\to 0}\int dk_z\frac{\alpha_s C_F\Lambda}{2\pi}\frac{[\delta(k_z-\bar x p_z)-\delta(\bar x p_z)]p_z}{k_z^2+\epsilon^2}
\eeq
其中 $\bar x=1-x$。另一方面，把式~\eqref{WLrenorm} 代入式~\eqref{qUnpolDef}，可以发现质量反项给出如下贡献
\begin{align}
&-\int \frac{dz}{2\pi}p_z\, e^{i(x-1)p_z z}|z|\, \delta m=-\lim_{\epsilon\to 0}\int \frac{dz}{2\pi}p_z\, e^{-i \bar x p_z z}\frac{1-e^{-\epsilon |z|}}{\epsilon}\delta m\non\\
&=-\lim_{\epsilon\to 0}\int \frac{dz}{2\pi}p_z\, e^{-i \bar x p_z z}\int_0^\infty \frac{d\alpha}{\sqrt{\pi \alpha}}(1-e^{-\frac{z^2}{4\alpha}})e^{-\alpha \epsilon^2}\delta m\non\\
&=-\lim_{\epsilon\to 0}\int \frac{dz}{2\pi}p_z\, e^{-i \bar x p_z z}\int_0^\infty d\alpha \int d k_z \frac{e^{-\alpha (k_z^2+\epsilon^2)}(1-e^{i k_z z})}{\pi} \delta m\non\\
&=-\lim_{\epsilon\to 0}\int \frac{dz}{2\pi}p_z\, e^{-i \bar x p_z z}\int dk_z \frac{1}{\pi}\frac{(1-e^{i k_z z})}{k_z^2+\epsilon^2}\delta m\non\\
&=-\lim_{\epsilon\to 0}\int \frac{dk_z}{\pi}p_z \frac{\delta(\bar x p_z)-\delta(k_z-\bar x p_z)}{k_z^2+\epsilon^2}\delta m.
\end{align}
因此我们得到
\beq
\delta m=-\frac{\alpha_s C_F}{2\pi}(\pi\Lambda)
\eeq
于单圈水平。

在坐标空间中，把威尔逊线展开到 $\mathcal O(g^2)$ 阶，我们有
\beq\label{WLcoordspace}
g^2 C_F\int_0^z dz_1\int_0^{z_1} dz_2 \frac{1}{4\pi^2}\frac{1}{(z_1-z_2)^2+a^2}=\frac{g^2 C_F}{4\pi^2}\big[\frac{z}{a}\tan^{-1}\big(\frac z a\big)-\frac 1 2\ln\big({1+\frac{z^2}{a^2}}\big)\big],
\eeq
其中我们引入截断 $a$，以正规化坐标空间胶子传播子在 $z_1\to z_2$ 处的短距离奇异性。经傅里叶变换到动量空间，上述结果导致
\beq\label{eq8}
\frac{g^2 C_F}{8\pi^2}\big[\frac{1}{a p_z(1-x)^2}-\frac{1}{|1-x|}+C\delta(1-x)\big],
\eeq
其中常数 $C$ 可由要求上式的 $x$ 积分为零来确定。于是我们可以认定
\beq
\Lambda=\frac{1}{a}.
\eeq

$\alpha_s$ 阶的质量反项图示于图~\ref{1loopFeynmassCT}。$\delta m$ 的规范无关性已可从下述事实看出：轴规范\B{~\cite{Xiong:2013bka}} 与费曼规范\B{~\cite{Ji:2015jwa}} 中的单圈修正包含完全相同的线性发散。利用 $R_\xi$ 规范容易发现，胶子传播子分子中额外的 $k^\mu k^\nu/k^2$ 项（$k$ 为胶子动量）对线性发散没有影响，因此 $\delta m$ 在不同规范中保持不变。


可以在图~\ref{1loopFeyn} 的单圈图上添加更多的胶子和夸克以构成更高圈图。幂次计数表明，只有那些以威尔逊线自能为子图的图才出现幂次发散。如附录所示，超出单圈之后，质量反项同样能消除威尔逊线的全部幂次发散。因此，计入质量反项贡献之后，准夸克分布中至多存在对数发散。这表明，就幂次发散而言，诸如准夸克分布这类双线性夸克算符的重正化与开放威尔逊线的重正化相同。在上述讨论中，我们证明了在微扰单圈水平上线性发散确实被质量反项抵消。在格点上，质量反项 $\delta m$ 必须非微扰地确定，例如可以按照文献~\cite{Musch:2010ka} 的方法进行。

\begin{figure}[tbp]
\centering
\includegraphics[width=0.2\textwidth]{images/1loopmassct}
\caption{费曼规范下准夸克分布的单圈质量反项图。}
\label{1loopFeynmassCT}
\end{figure}
